\phantomsection 
\part{線型空間}

\begin{abstract}
  この章では，線型空間を導入し，ベクトルを抽象的な立場から捉える．
  これにより，数ベクトルだけでなく，行列や関数なども同じ枠組みで扱えるようになる．
  さらに，線型部分空間，基底，次元などの基本概念を導入し，線型空間の構造について考察する．
\end{abstract}

\phantomsection
\section{線型空間}

\phantomsection
\subsection{線型空間の定義}


線型代数学では，ある条件を満たす集合の元を「ベクトル」として扱う．
そこで，その枠組みとなる\textbf{線型空間}\index{せんけいくうかん@線型空間}を定義する．

\partref{行列と連立一次方程式}から\partref{行列式}までは$K=\mathbb{R}$または$K=\mathbb{C}$としていたが，
本章では$K$を任意の\textbf{体}\index{たい@体}（field）とする%
\footnote{体とは，大まかにいえば，加法・減法・乗法，および$0$でない元による除法が自由に行える代数系である．
$\mathbb{R}$や$\mathbb{C}$はその代表例である．}．

次の定義では，体$K$の加法・乗法と
$V$の加法・スカラー倍を区別するため，それぞれ異なる記号を用いる．

\begin{definition}{線型空間}{線型空間}
  $K$を体とし，その加法と乗法をそれぞれ
  \[
    \Kadd\colon K\times K\to K,
    \qquad
    \Kmul\colon K\times K\to K
  \]
  とする．また，加法の零元を$0_K$，乗法の単位元を$1_K$とする．

  集合$V$に，加法とスカラー倍
  \[
    \Vadd\colon V\times V\to V,
    \qquad
    \scalmul\colon K\times V\to V
  \]
  が定められているとする．
  このとき，次の条件を満たす$V$を体$K$上の\textbf{線型空間}という．

  \begin{enumerate}[label=(\roman*),leftmargin=3.2em]
    \item \textbf{加法に関して可換群}

    $V$は，加法$\Vadd$に関して可換群をなす．

    \item \textbf{スカラー倍に関する性質}

    任意の$\bm{v},\bm{w}\in V$および$\alpha,\beta\in K$に対して，
    \[
      1_K\scalmul\bm{v}=\bm{v},
      \qquad
      (\alpha\Kadd\beta)\scalmul\bm{v}
      =
      (\alpha\scalmul\bm{v})\Vadd(\beta\scalmul\bm{v}),
    \]
    \[
      \alpha\scalmul(\bm{v}\Vadd\bm{w})
      =
      (\alpha\scalmul\bm{v})\Vadd(\alpha\scalmul\bm{w}),
      \qquad
      (\alpha\Kmul\beta)\scalmul\bm{v}
      =
      \alpha\scalmul(\beta\scalmul\bm{v})
    \]
    が成り立つ．
  \end{enumerate}

  $V$の元を\textbf{ベクトル}という．
\end{definition}

この定義では，$4$つの演算を区別するために，
一時的に$\Kadd$，$\Kmul$，$\Vadd$，$\scalmul$という記号を用いた．
しかし，どの演算であるかは文脈から判断できるため，以下では
\[
  \alpha\Kadd\beta,\quad
  \bm{v}\Vadd\bm{w},\quad
  \alpha\Kmul\beta,\quad
  \alpha\scalmul\bm{v}
\]
を，それぞれ
\[
  \alpha+\beta,\quad
  \bm{v}+\bm{w},\quad
  \alpha\beta,\quad
  \alpha\bm{v}
\]
と表す．
たとえば
\[
  (\alpha+\beta)\bm{v}
  =
  \alpha\bm{v}+\beta\bm{v}
\]
では，左辺の$+$は$K$の加法，右辺の$+$は$V$の加法である．

\begin{remark}[加法に関する可換群の条件]
線型空間の公理のうち，
「加法に関して可換群である」という条件には，

\begin{itemize}
  \item 結合法則
  \item 交換法則
  \item 零ベクトルの存在
  \item 加法逆元の存在
\end{itemize}

が含まれている．
具体的には，任意の$\bm{v},\bm{w},\bm{u}\in V$に対して，
\[
  (\bm{v}+\bm{w})+\bm{u}
  =
  \bm{v}+(\bm{w}+\bm{u}),
\]
\[
  \bm{v}+\bm{w}
  =
  \bm{w}+\bm{v},
\]
\[
  \bm{v}+\bm{0}_V
  =
  \bm{v},
\]
\[
  \bm{v}+(-\bm{v})
  =
  \bm{0}_V
\]
が成り立つ．
ここではすでに通常の表記に戻しているため，
$V$の加法には通常の$+$を用いている．
\end{remark}

\begin{example}{線型空間の例}{線型空間の例}
次の集合はいずれも，通常の加法とスカラー倍に関して線型空間となる．

\begin{itemize}
\item $\mathbb{R}$上の線型空間$\mathbb{R}^n$
\item $\mathbb{C}$上の線型空間$\mathbb{C}^n$
\item $K$上の線型空間$M_{m,n}(K)$（$K$の元を成分とする$m\times n$行列全体，\deref{行列全体の集合}）
\item $\mathbb{R}$上の線型空間$\mathbb{R}[x]$（実数係数多項式全体）
\end{itemize}

これらは一見異なる対象であるが，
線型空間という共通の枠組みで扱うことができる\footnote{同じ集合でも，どの体の上の線型空間とみるかによって構造は異なる．たとえば$\mathbb{C}^n$は，スカラーを実数に制限すれば$\mathbb{R}$上の線型空間ともみなせる．線型空間を考えるときに係数体を明示するのはこのためである．}．
\end{example}

線型空間の公理からは，数ベクトルの場合と同様の計算規則が導かれる．

\begin{prop}{線型空間の基本性質}{線型空間の基本性質}
  $V$を体$K$上の線型空間とする．
  任意の$\bm{v} \in V$および任意の$\alpha \in K$に対して，次が成り立つ：
  \begin{enumerate}[label=(\roman*),leftmargin=3.2em]
    \item $0_K \bm{v} = \bm{0}_V$
    \item $\alpha \bm{0}_V = \bm{0}_V$
    \item $(-1_K)\bm{v} = -\bm{v}$
  \end{enumerate}
\end{prop}

\begin{leftbar}
\begin{proof}
  (i)について，スカラー倍に関する分配律より，
  \[
    0_K \bm{v} = (0_K + 0_K)\bm{v} = 0_K \bm{v} + 0_K \bm{v}
  \]
  である．両辺に$0_K \bm{v}$の加法逆元$-(0_K \bm{v})$を加えると，
  \[
    \bm{0}_V = 0_K \bm{v}
  \]
  を得る．

  (ii)について，同様に，
  \[
    \alpha \bm{0}_V = \alpha(\bm{0}_V + \bm{0}_V) = \alpha \bm{0}_V + \alpha \bm{0}_V
  \]
  であるから，両辺に$-(\alpha \bm{0}_V)$を加えて$\bm{0}_V = \alpha \bm{0}_V$を得る．

  (iii)について，(i)より，
  \[
    \bm{v} + (-1_K)\bm{v} = 1_K \bm{v} + (-1_K)\bm{v} = \bigl(1_K + (-1_K)\bigr)\bm{v} = 0_K \bm{v} = \bm{0}_V
  \]
  である．加法逆元の一意性より，$(-1_K)\bm{v} = -\bm{v}$である．
  以上の考察により，命題が示された．
\end{proof}
\end{leftbar}

\subsection{線型部分空間}

線型空間の部分集合のうち，
それ自身も線型空間となるものを線型部分空間という．

\begin{definition}{線型部分空間}{線型部分空間}
  体$K$上の線型空間$V$の部分集合$W$が，
  次の条件を満たすとき，
  $W$を$V$の\textbf{線型部分空間}という．
\begin{enumerate}[label=(\roman*),leftmargin=3.2em]
    \item \textbf{空でない}
    \[
      W\neq\varnothing
    \]
    である．

    \item \textbf{加法に関する性質}

    任意の$\bm{v},\bm{w}\in W$に対して，
    \[
      \bm{v}+\bm{w}\in W
    \]
    である．

    \item \textbf{スカラー倍に関する性質}

    任意の$\bm{v}\in W$，
    任意の$\alpha\in K$に対して，
    \[
      \alpha \bm{v}\in W
    \]
    である．
  \end{enumerate}
\end{definition}

\begin{remark}[空でないことと零ベクトル]
定義の最初の条件は
\[
\bm{0}_V\in W
\]
としてもよい．

実際，$W$が空でなければ$\bm{w} \in W$がとれ，$W$がスカラー倍について閉じていることと\prref{線型空間の基本性質}より，
\[
\bm{0}_V = 0_K \bm{w} \in W
\]
となる．
\end{remark}

線型部分空間の例をいくつか見よう．

\begin{example}{線型部分空間の例}{線型部分空間の例}
  \begin{enumerate}[label=(\arabic*),leftmargin=3.2em]
    \item 任意の線型空間$V$に対して，$\{\bm{0}_V\}$と$V$自身は，いずれも$V$の線型部分空間である．
    \item $K = \mathbb{R}$とする．原点を通る直線
      \[
        W = \left\{ \begin{pmatrix} x \\ y \end{pmatrix} \in \mathbb{R}^2 ~\middle|~ 2x - y = 0 \right\}
      \]
      は$\mathbb{R}^2$の線型部分空間である．
      実際，$\bm{0} \in W$であり，$2x_1 - y_1 = 0$かつ$2x_2 - y_2 = 0$ならば$2(x_1 + x_2) - (y_1 + y_2) = 0$，また任意の$\alpha \in \mathbb{R}$に対して$2(\alpha x_1) - \alpha y_1 = \alpha(2x_1 - y_1) = 0$となるからである．
      一方，原点を通らない直線（たとえば$2x - y = 1$で定まる直線）は$\bm{0}$を含まないので，線型部分空間ではない．
    \item 次数が$n$以下の実数係数多項式全体（零多項式を含む）は，$\mathbb{R}[x]$の線型部分空間である．
      次数$n$以下の多項式どうしの和や，そのスカラー倍によって，次数が$n$を超えることはないからである．
  \end{enumerate}
\end{example}

さらに，\partref{線型写像}で導入した線型写像の像と核（\deref{像と核}）は，線型部分空間の重要な例を与える\footnote{\partref{線型写像}では$K$を$\mathbb{R}$または$\mathbb{C}$としていたが，線型写像とその像・核の定義（\deref{線型写像}，\deref{像と核}），および\prref{線型写像と零ベクトル}や\prref{線型写像と線型結合}などの基本性質は，一般の体$K$に対してもそのまま通用する．}．

\begin{prop}{像と核の線型部分空間性}{像と核の線型部分空間性}
  $f \colon K^n \to K^m$を線型写像とする．このとき，次が成り立つ．
  \begin{enumerate}[label=(\roman*),leftmargin=3.2em]
    \item $\Ker f$は$K^n$の線型部分空間である．
    \item $\Im f$は$K^m$の線型部分空間である．
  \end{enumerate}
\end{prop}

\begin{leftbar}
\begin{proof}
  (i)について示す．
  \prref{線型写像と零ベクトル}より$f(\bm{0}) = \bm{0}$であるから，$\bm{0} \in \Ker f$であり，$\Ker f$は空でない．
  $\bm{x}, \bm{y} \in \Ker f$とすると，$f$の線型性より
  \[
    f(\bm{x} + \bm{y}) = f(\bm{x}) + f(\bm{y}) = \bm{0} + \bm{0} = \bm{0}
  \]
  であるから，$\bm{x} + \bm{y} \in \Ker f$である．
  また，任意の$\alpha \in K$に対して
  \[
    f(\alpha\bm{x}) = \alpha f(\bm{x}) = \alpha\bm{0} = \bm{0}
  \]
  であるから，$\alpha\bm{x} \in \Ker f$である．
  よって$\Ker f$は$K^n$の線型部分空間である．

  (ii)について示す．
  $\bm{0} = f(\bm{0}) \in \Im f$であるから，$\Im f$は空でない．
  $\bm{v}, \bm{w} \in \Im f$とすると，$\bm{v} = f(\bm{x})$，$\bm{w} = f(\bm{y})$となる$\bm{x}, \bm{y} \in K^n$がとれる．$f$の線型性より
  \[
    \bm{v} + \bm{w} = f(\bm{x}) + f(\bm{y}) = f(\bm{x} + \bm{y}) \in \Im f
  \]
  であり，また任意の$\alpha \in K$に対して
  \[
    \alpha\bm{v} = \alpha f(\bm{x}) = f(\alpha\bm{x}) \in \Im f
  \]
  である．よって$\Im f$は$K^m$の線型部分空間である．
  これが証明すべきことであった．
\end{proof}
\end{leftbar}

とくに，斉次連立一次方程式$A\bm{x} = \bm{0}$の解全体は$\Ker T_A$であったから（\prref{連立一次方程式と像と核}），$K^n$の線型部分空間である．
これを$A\bm{x} = \bm{0}$の\textbf{解空間}\index{かいくうかん@解空間}という．
「$2$つの解の和はふたたび解であり，解のスカラー倍もまた解である」という重ね合わせの性質こそが，部分空間性の意味するところである．

線型部分空間どうしから新たな線型部分空間を作る方法として，
共通部分と和が重要である．

\begin{definition}{共通部分・和}{共通部分・和}
  $V$を線型空間，
  $W_1,W_2$を$V$の線型部分空間とする．

  このとき，
  \[
  W_1\cap W_2
  \]
  を$W_1$と$W_2$の\textbf{共通部分}という．

  また，
  \[
  W_1+W_2
  =
  \{
  \bm{w}_1+\bm{w}_2
  \mid
  \bm{w}_1\in W_1,\
  \bm{w}_2\in W_2
  \}
  \]
  を$W_1$と$W_2$の\textbf{和}という．
\end{definition}

\begin{prop}{共通部分と和の線型部分空間性}{共通部分と和の線型部分空間性}
  \deref{共通部分・和}で定義した，
  線型部分空間の共通部分および和は，
  ともに線型部分空間である．
\end{prop}

\begin{leftbar}
\begin{proof}
  まず，共通部分$W_1\cap W_2$について示す．
  $W_1,W_2$はいずれも線型部分空間であるから，
  \[
    \bm{0}_V\in W_1,
    \quad
    \bm{0}_V\in W_2
  \]
  である．したがって，
  \[
    \bm{0}_V\in W_1\cap W_2
  \]
  であり，$W_1\cap W_2$は空でない．

  次に，$\bm{x},\bm{y}\in W_1\cap W_2$とする．
  このとき，
  \[
    (\bm{x},\bm{y}\in W_1)
    \quad \land \quad
    (\bm{x},\bm{y}\in W_2)
  \]
  である．$W_1,W_2$は加法について閉じているから，
  \[
    (\bm{x}+\bm{y}\in W_1)
    \quad \land \quad
    (\bm{x}+\bm{y}\in W_2)
  \]
  となる．ゆえに，
  \[
    \bm{x}+\bm{y}\in W_1\cap W_2
  \]
  である．

  また，任意の$\alpha\in K$に対して，
  $W_1,W_2$はスカラー倍について閉じているから，
  \[
    (\alpha\bm{x}\in W_1)
    \quad \land \quad
    (\alpha\bm{x}\in W_2)
  \]
  となる．したがって，
  \[
    \alpha\bm{x}\in W_1\cap W_2
  \]
  である．以上より，$W_1\cap W_2$は$V$の線型部分空間である．

  続いて，和$W_1+W_2$について示す．
  \[
    \bm{0}_V=\bm{0}_V+\bm{0}_V
  \]
  であり，$\bm{0}_V\in W_1$かつ$\bm{0}_V\in W_2$であるから，
  \[
    \bm{0}_V\in W_1+W_2
  \]
  である．したがって，$W_1+W_2$は空でない．

  $\bm{x},\bm{y}\in W_1+W_2$とする．
  このとき，ある$\bm{x}_1,\bm{y}_1\in W_1$および
  $\bm{x}_2,\bm{y}_2\in W_2$が存在して，
  \[
    \bm{x}=\bm{x}_1+\bm{x}_2,
    \quad
    \bm{y}=\bm{y}_1+\bm{y}_2
  \]
  と表せる．したがって，
  \begin{align*}
    \bm{x}+\bm{y}
    &=(\bm{x}_1+\bm{x}_2)+(\bm{y}_1+\bm{y}_2)\\
    &=(\bm{x}_1+\bm{y}_1)+(\bm{x}_2+\bm{y}_2)
  \end{align*}
  である．ここで，
  \[
    \bm{x}_1+\bm{y}_1\in W_1,
    \quad
    \bm{x}_2+\bm{y}_2\in W_2
  \]
  であるから，
  \[
    \bm{x}+\bm{y}\in W_1+W_2
  \]
  となる．

  さらに，任意の$\alpha\in K$に対して，
  \[
    \alpha\bm{x}
    =\alpha(\bm{x}_1+\bm{x}_2)
    =\alpha\bm{x}_1+\alpha\bm{x}_2
  \]
  である．$\alpha\bm{x}_1\in W_1$かつ
  $\alpha\bm{x}_2\in W_2$であるから，
  \[
    \alpha\bm{x}\in W_1+W_2
  \]
  である．以上より，$W_1+W_2$も$V$の線型部分空間である．
\end{proof}
\end{leftbar}

\phantomsection
\subsection{線型結合と線型包}

\partref{線型独立性と階数}では，数ベクトルに対して線型結合を定義した（\deref{線型結合}）．
その定義は，一般の線型空間の元に対しても，まったく同様に通用する．

\begin{definition}{線型結合（一般の線型空間）}{一般の線型結合}
  $V$を体$K$上の線型空間とする．
  $n$個のベクトル$\bm{v}_1, \bm{v}_2, \dots, \bm{v}_n \in V$に対して，
  \[
    \lambda_1 \bm{v}_1 + \lambda_2 \bm{v}_2 + \cdots + \lambda_n \bm{v}_n
    \quad (\lambda_i \in K)
  \]
  の形のベクトルを，$\bm{v}_1, \bm{v}_2, \dots, \bm{v}_n$の\textbf{線型結合}という．

  線型関係，線型独立，線型従属についても，\partref{線型独立性と階数}と同じ式によって定義される．
\end{definition}

ベクトルの集合が与えられたとき，その元の線型結合として表されるベクトルをすべて集めた集合を考えよう．

\begin{definition}{線型包}{線型包}
  $V$を体$K$上の線型空間とし，$S$を$V$の部分集合とする．

  $S \neq \varnothing$のとき，$S$の有限個の元の線型結合全体の集合
  \[
    \Span(S)
    = \{ \lambda_1 \bm{v}_1 + \lambda_2 \bm{v}_2 + \cdots + \lambda_n \bm{v}_n \mid n \ge 1,\ \bm{v}_1, \dots, \bm{v}_n \in S,\ \lambda_1, \dots, \lambda_n \in K \}
  \]
  を$S$の\textbf{線型包}\index{せんけいほう@線型包}\index{Span@$\Span$}（\textbf{張る空間}\index{はるくうかん@張る空間}，生成される部分空間）という．
  また，$\Span(\varnothing) = \{\bm{0}_V\}$と定める．

  有限個のベクトル$\bm{v}_1, \bm{v}_2, \dots, \bm{v}_n \in V$に対しては，
  \[
    \Span(\bm{v}_1, \bm{v}_2, \dots, \bm{v}_n) = \Span(\{\bm{v}_1, \bm{v}_2, \dots, \bm{v}_n\})
  \]
  ともかく．
\end{definition}

\begin{example}{線型包の例}{線型包の例}
  $K^3$の標準基底ベクトル$\bm{e}_1$，$\bm{e}_2$に対して，
  \[
    \Span(\bm{e}_1, \bm{e}_2)
    = \left\{ \lambda_1 \begin{pmatrix} 1 \\ 0 \\ 0 \end{pmatrix} + \lambda_2 \begin{pmatrix} 0 \\ 1 \\ 0 \end{pmatrix} ~\middle|~ \lambda_1, \lambda_2 \in K \right\}
    = \left\{ \begin{pmatrix} \lambda_1 \\ \lambda_2 \\ 0 \end{pmatrix} ~\middle|~ \lambda_1, \lambda_2 \in K \right\}
  \]
  である．$K = \mathbb{R}$のとき，これは$xyz$空間における$xy$平面にほかならない．

  また，線型空間$V$の$1$個のベクトル$\bm{v} \neq \bm{0}_V$に対して，
  \[
    \Span(\bm{v}) = \{ \lambda \bm{v} \mid \lambda \in K \}
  \]
  であり，たとえば$V = \mathbb{R}^3$のとき，これは原点を通り$\bm{v}$の方向に伸びる直線である．

  さらに，$V$を実数係数多項式全体のなす線型空間とすると，
  \[
    \Span(1, x, x^2)
    = \{ a + bx + cx^2 \mid a, b, c \in \mathbb{R} \}
  \]
  であり，これは次数が$2$以下の多項式全体（零多項式を含む）にほかならない．
  このように，線型包の概念は数ベクトルに限らず適用できる．
\end{example}

線型包は，つねに線型部分空間となる．

\begin{prop}{線型包の線型部分空間性}{線型包の線型部分空間性}
  $V$を体$K$上の線型空間，$S$を$V$の部分集合とする．
  このとき，$\Span(S)$は$V$の線型部分空間である．
\end{prop}

\begin{leftbar}
\begin{proof}
  まず，$S = \varnothing$のときを考える．
  このとき$\Span(\varnothing) = \{\bm{0}_V\}$であり，
  \[
    \bm{0}_V + \bm{0}_V = \bm{0}_V
  \]
  および，任意の$\alpha \in K$に対して\prref{線型空間の基本性質}より
  \[
    \alpha \bm{0}_V = \bm{0}_V
  \]
  が成り立つ．よって$\{\bm{0}_V\}$は空でなく，加法とスカラー倍について閉じているから，$V$の線型部分空間である．

  次に，$S \neq \varnothing$とする．
  $\bm{v} \in S$をとると，
  \[
    \bm{v} = 1_K \bm{v} \in \Span(S)
  \]
  であるから，$\Span(S)$は空でない．

  $\bm{x}, \bm{y} \in \Span(S)$とすると，ある$\bm{v}_1, \dots, \bm{v}_m \in S$，$\bm{w}_1, \dots, \bm{w}_l \in S$および$\lambda_i, \mu_j \in K$を用いて
  \[
    \bm{x} = \sum_{i=1}^{m} \lambda_i \bm{v}_i, \quad
    \bm{y} = \sum_{j=1}^{l} \mu_j \bm{w}_j
  \]
  と表せる．このとき，
  \[
    \bm{x} + \bm{y}
    = \lambda_1 \bm{v}_1 + \cdots + \lambda_m \bm{v}_m + \mu_1 \bm{w}_1 + \cdots + \mu_l \bm{w}_l
  \]
  も$S$の有限個の元の線型結合であるから，$\bm{x} + \bm{y} \in \Span(S)$である．

  さらに，任意の$\alpha \in K$に対して，スカラー倍に関する分配律より，
  \[
    \alpha\bm{x} = \sum_{i=1}^{m} (\alpha\lambda_i) \bm{v}_i
  \]
  も$S$の有限個の元の線型結合であるから，$\alpha\bm{x} \in \Span(S)$である．
  以上より，$\Span(S)$は$V$の線型部分空間である．
\end{proof}
\end{leftbar}

さらに，線型包は「$S$を含む最小の線型部分空間」として特徴づけられる．

\begin{prop}{線型包の最小性}{線型包の最小性}
  $V$を体$K$上の線型空間，$S$を$V$の部分集合とする．このとき，次が成り立つ：
  \begin{enumerate}[label=(\roman*),leftmargin=3.2em]
    \item $S \subset \Span(S)$である．
    \item $W$が$S \subset W$を満たす$V$の線型部分空間ならば，$\Span(S) \subset W$である．
  \end{enumerate}
  すなわち，$\Span(S)$は，$S$を含む最小の線型部分空間である．
\end{prop}

\begin{leftbar}
\begin{proof}
  (i)について，任意の$\bm{v} \in S$は
  \[
    \bm{v} = 1_K \bm{v}
  \]
  と表せるから，$\bm{v} \in \Span(S)$である．

  (ii)について，まず$S = \varnothing$のときを考える．
  任意の線型部分空間$W$は必ず零ベクトルを元にもつ（\deref{線型部分空間}のあとのRemarkを参照）から，
  \[
    \Span(\varnothing) = \{\bm{0}_V\} \subset W
  \]
  である．

  次に，$S \neq \varnothing$とし，$\Span(S)$の任意の元
  \[
    \bm{x} = \lambda_1 \bm{v}_1 + \lambda_2 \bm{v}_2 + \cdots + \lambda_n \bm{v}_n
    \quad (\bm{v}_i \in S,\ \lambda_i \in K)
  \]
  をとる．$S \subset W$より$\bm{v}_i \in W$であり，$W$はスカラー倍について閉じているから，
  \[
    \lambda_i \bm{v}_i \in W \quad (1 \le i \le n)
  \]
  である．さらに，$W$は加法について閉じているから，これらの和である$\bm{x}$も$W$に属する．
  よって$\Span(S) \subset W$である．
  以上の考察により，命題が示された．
\end{proof}
\end{leftbar}

\begin{remark}[線型包と共通部分]
\prref{線型包の最小性}より，$\Span(S)$は「$S$を含む線型部分空間すべての共通部分」とも一致する．

実際，\prref{線型包の線型部分空間性}と\prref{線型包の最小性}(i)より，$\Span(S)$自身が「$S$を含む線型部分空間」のひとつであるから，共通部分は$\Span(S)$に含まれる．
逆に，\prref{線型包の最小性}(ii)より，$\Span(S)$は$S$を含むどの線型部分空間にも含まれるから，共通部分にも含まれる．

なお，\prref{共通部分と和の線型部分空間性}では二つの線型部分空間の共通部分を扱ったが，同様の議論により，任意個の線型部分空間の共通部分も線型部分空間となる．
\end{remark}

和と線型包のあいだには，次の関係が成り立つ．

\begin{prop}{和と線型包}{和と線型包}
  $V$を線型空間，$W_1, W_2$を$V$の線型部分空間とする．このとき，
  \[
    W_1 + W_2 = \Span(W_1 \cup W_2)
  \]
  が成り立つ．
\end{prop}

\begin{leftbar}
\begin{proof}
  まず，$W_1 + W_2 \subset \Span(W_1 \cup W_2)$を示す．
  任意の$\bm{w}_1 \in W_1$，$\bm{w}_2 \in W_2$に対して，
  \[
    \bm{w}_1 + \bm{w}_2 = 1_K \bm{w}_1 + 1_K \bm{w}_2
  \]
  は$W_1 \cup W_2$の元の線型結合であるから，$\bm{w}_1 + \bm{w}_2 \in \Span(W_1 \cup W_2)$である．

  逆に，$\Span(W_1 \cup W_2) \subset W_1 + W_2$を示す．
  \prref{共通部分と和の線型部分空間性}より$W_1 + W_2$は$V$の線型部分空間であり，
  \[
    \bm{w}_1 = \bm{w}_1 + \bm{0}_V \in W_1 + W_2,
    \quad
    \bm{w}_2 = \bm{0}_V + \bm{w}_2 \in W_1 + W_2
  \]
  より$W_1 \cup W_2 \subset W_1 + W_2$である．
  よって，\prref{線型包の最小性}より$\Span(W_1 \cup W_2) \subset W_1 + W_2$である．
  以上の考察により，命題が示された．
\end{proof}
\end{leftbar}

\begin{remark}[合併は線型部分空間とは限らない]
一般に，線型部分空間どうしの合併$W_1 \cup W_2$は線型部分空間になるとは限らない．
たとえば，$\mathbb{R}^2$において$x$軸と$y$軸はいずれも線型部分空間であるが，
\[
  \begin{pmatrix} 1 \\ 0 \end{pmatrix}
  +
  \begin{pmatrix} 0 \\ 1 \end{pmatrix}
  =
  \begin{pmatrix} 1 \\ 1 \end{pmatrix}
\]
は$x$軸にも$y$軸にも属さないので，合併は加法について閉じていない．

\prref{和と線型包}は，合併そのものではなく，和$W_1 + W_2$こそが「$W_1$と$W_2$の両方を含む最小の線型部分空間」であることを示している．
\end{remark}

\phantomsection
\subsection{直和}

本項では「直和」とよばれる二つの構成を扱う．
まず，二つの線型空間を組にして一つの線型空間を作る方法として，外部直和を定義する．

\begin{definition}{直和（外部直和）}{直和}
  $V,W$を体$K$上の線型空間とする．

  直積集合
  \[
  V\times W
  =
  \{
  (\bm{v},\bm{w})
  \mid
  \bm{v}\in V,\
  \bm{w}\in W
  \}
  \]
  を考え，
  加法およびスカラー倍を
  \[
  (\bm{v}_1,\bm{w}_1)+(\bm{v}_2,\bm{w}_2)
  =
  (\bm{v}_1+\bm{v}_2,\bm{w}_1+\bm{w}_2),
  \]
  \[
  \alpha(\bm{v},\bm{w})
  =
  (\alpha \bm{v},\alpha \bm{w})
  \]
  によって定める．

  このとき，線型空間の各公理を成分ごとに確認することにより，
  $V\times W$はこれらの演算に関して体$K$上の線型空間となる．

  これを
  $V$と$W$の\textbf{直和}\index{ちょくわ@直和}（\textbf{外部直和}\index{がいぶちょくわ@外部直和}）といい，
  \[
  V\oplus W
  \]
  と表す．
\end{definition}

\begin{remark}[外部直和の直観]
直和
$V\oplus W$
は，
二つの線型空間を互いに独立な成分として並べた線型空間と考えることができる．
\end{remark}

一方，一つの線型空間$V$の「内部」で，二つの線型部分空間が$V$を過不足なく分担している状況を表す概念として，内部直和を定義する．

\begin{definition}{内部直和}{内部直和}
  $W_1, W_2$を体$K$上の線型空間$V$の線型部分空間とする．
  \[
    V = W_1 + W_2,
    \qquad
    W_1 \cap W_2 = \{\bm{0}_V\}
  \]
  がともに成り立つとき（和$W_1 + W_2$については\deref{共通部分・和}を参照），$V$は$W_1$と$W_2$の\textbf{内部直和}\index{ないぶちょくわ@内部直和}であるといい，
  \[
    V = W_1 \oplus W_2
  \]
  と表す．
\end{definition}

内部直和は，「ベクトルの分解の一意性」によって特徴づけられる．

\begin{prop}{内部直和と分解の一意性}{内部直和と分解の一意性}
  $W_1, W_2$を線型空間$V$の線型部分空間とする．このとき，次の二つは同値である：
  \begin{enumerate}[label=(\roman*),leftmargin=3.2em]
    \item $V = W_1 \oplus W_2$（内部直和）である．
    \item 任意の$\bm{v} \in V$は，
    \[
      \bm{v} = \bm{w}_1 + \bm{w}_2
      \quad (\bm{w}_1 \in W_1,~ \bm{w}_2 \in W_2)
    \]
    の形にただ一通りに表される．
  \end{enumerate}
\end{prop}

\begin{leftbar}
\begin{proof}
  まず，(i)を仮定して(ii)を示す．
  $V = W_1 + W_2$であるから，任意の$\bm{v} \in V$は$\bm{v} = \bm{w}_1 + \bm{w}_2$（$\bm{w}_1 \in W_1$，$\bm{w}_2 \in W_2$）の形に表せる．
  いま，
  \[
    \bm{v} = \bm{w}_1 + \bm{w}_2 = \bm{w}'_1 + \bm{w}'_2
    \quad (\bm{w}_i, \bm{w}'_i \in W_i)
  \]
  と二通りに表せたとすると，移項して
  \[
    \bm{w}_1 - \bm{w}'_1 = \bm{w}'_2 - \bm{w}_2
  \]
  を得る．左辺は$W_1$に，右辺は$W_2$に属するから，両辺が表すベクトルは$W_1 \cap W_2 = \{\bm{0}_V\}$に属する．
  よって$\bm{w}_1 = \bm{w}'_1$かつ$\bm{w}_2 = \bm{w}'_2$であり，表し方はただ一通りである．

  逆に，(ii)を仮定して(i)を示す．
  任意の$\bm{v} \in V$が$\bm{w}_1 + \bm{w}_2$の形に表せることから，$V = W_1 + W_2$である．
  次に，$\bm{v} \in W_1 \cap W_2$をとると，
  \[
    \bm{v} = \bm{v} + \bm{0}_V = \bm{0}_V + \bm{v}
  \]
  は，いずれも「$W_1$の元と$W_2$の元の和」としての$\bm{v}$の表示である．
  表し方の一意性より$\bm{v} = \bm{0}_V$となるから，$W_1 \cap W_2 = \{\bm{0}_V\}$である．
  以上の考察により，命題が示された．
\end{proof}
\end{leftbar}

\begin{remark}[内部直和と外部直和の同一視]
  外部直和と内部直和は異なる構成であるが，同じ記号$\oplus$で表される．
  これは，$V = W_1 \oplus W_2$（内部直和）のとき，外部直和$W_1 \oplus W_2 = W_1 \times W_2$の元$(\bm{w}_1, \bm{w}_2)$に$\bm{w}_1 + \bm{w}_2 \in V$を対応させる写像が全単射となり（全射性は$V = W_1 + W_2$から，単射性は\prref{内部直和と分解の一意性}の分解の一意性から従う），さらに和とスカラー倍も保たれるため，両者を自然に同一視できるからである．
\end{remark}

\phantomsection
\section{基底と次元}

\phantomsection
\subsection{基底}

前節で，ベクトルの集合$S$から生成される線型部分空間$\Span(S)$を導入した．
一方，\partref{線型独立性と階数}では，ベクトルの組に「無駄がない」ことを表す概念として線型独立を導入した（\deref{線型独立}）．

この二つを組み合わせよう．
すなわち，「$V$全体を生成し，かつ無駄がない」ベクトルの組を考える．

\begin{definition}{基底}{基底}
  $V$を体$K$上の線型空間とする．
  $V$のベクトルの組$\mathcal{B} = (\bm{v}_1, \bm{v}_2, \dots, \bm{v}_n)$が次の二つの条件を満たすとき，
  $\mathcal{B}$を$V$の\textbf{基底}\index{きてい@基底}という．
  \begin{enumerate}[label=(\roman*),leftmargin=3.2em]
    \item \textbf{線型独立性}

    $\bm{v}_1, \bm{v}_2, \dots, \bm{v}_n$は線型独立である．

    \item \textbf{生成性}
    \[
      \Span(\bm{v}_1, \bm{v}_2, \dots, \bm{v}_n) = V
    \]
    である．
  \end{enumerate}

  また，$V = \{\bm{0}_V\}$のときは，空の組$(\ )$を$V$の基底と定める．
\end{definition}

\begin{remark}[基底を「組」として定義する理由]
  基底を「集合」ではなく「組」（順序のついた有限列）として定義していることに注意してほしい．
  順序を込めて考えるのは，のちに定義する座標（\deref{座標ベクトル}）が「第何成分か」を指定して初めて意味をもつからである．
  文献によっては基底を集合として定義し，座標を考えるときにあらためて順序を固定する流儀もある．

  また，$V = \{\bm{0}_V\}$に対して空の組を基底と定めたのは，$\Span(\varnothing) = \{\bm{0}_V\}$という約束（\deref{線型包}）と整合させるためである．
  「$\bm{0}_V$のみからなる組」を基底とすることはできない．
  実際，$1_K \bm{0}_V = \bm{0}_V$は非自明な線型関係であり，$(\bm{0}_V)$は線型従属だからである．
\end{remark}

\begin{example}{基底の例}{基底の例}
  \begin{enumerate}[label=(\arabic*),leftmargin=3.2em]
    \item $K^n$において，標準基底ベクトルの組$(\bm{e}_1, \bm{e}_2, \dots, \bm{e}_n)$は基底である．
    これを$K^n$の\textbf{標準基底}\index{ひょうじゅんきてい@標準基底}という．

    \item $K^2$において，
    \[
      \bm{v}_1 = \begin{pmatrix} 1 \\ 1 \end{pmatrix}, \quad
      \bm{v}_2 = \begin{pmatrix} 1 \\ -1 \end{pmatrix}
    \]
    の組も基底である．このように，基底の取り方は一通りではない．

    \item $K$の元を成分とする$m \times n$行列全体のなす線型空間$M_{m,n}(K)$において，$(i,j)$成分のみが$1$で他が$0$である行列$E_{ij}$の組は基底である\footnote{この$E_{ij}$は\textbf{行列単位}とよばれる．単位行列$E_n$（\deref{零行列と単位行列}）と記号が似ているが別物である．なお，正方行列（$m = n$）の場合には$E_n = E_{11} + E_{22} + \cdots + E_{nn}$という関係がある．}．

    \item 次数が$n$以下の実数係数多項式全体のなす線型空間において，$(1, x, x^2, \dots, x^n)$は基底である．
  \end{enumerate}
\end{example}

\begin{definition}{有限生成}{有限生成}
  線型空間$V$に対して，有限個のベクトル$\bm{v}_1, \bm{v}_2, \dots, \bm{v}_n \in V$が存在して
  \[
    V = \Span(\bm{v}_1, \bm{v}_2, \dots, \bm{v}_n)
  \]
  となるとき，$V$は\textbf{有限生成}\index{ゆうげんせいせい@有限生成}であるという．
\end{definition}

\begin{mycolumn}
  有限生成でない線型空間も存在する．
  たとえば，実数係数多項式全体のなす線型空間$\mathbb{R}[x]$は有限生成ではない．
  実際，有限個の多項式$f_1, \dots, f_n$に対して，十分大きい整数$d$をとれば，$f_1, \dots, f_n$はいずれも次数$d$以下の多項式である．次数$d$以下の多項式（零多項式を含む）の線型結合はふたたび次数$d$以下の多項式であるから，たとえば$x^{d+1}$は$f_1, \dots, f_n$の線型結合として表せず，より高次の多項式は生成できない．

  このような空間にも「基底」を考えることはできるが，そこでは無限個のベクトルからなる集合を扱う必要があり，線型独立性も「その有限部分集合がつねに線型独立である」と定義し直さなければならない．
  さらに，一般の線型空間に基底が存在することの証明には選択公理（ツォルンの補題）が必要となる．

  本書では，無限次元の線型空間を本格的に扱うことはせず，有限生成性（のちの言葉でいえば有限次元性）を仮定する命題では，その都度これを明記する．
  また，本書で単に「基底」というときは，つねに\deref{基底}の意味での，有限個のベクトルからなる順序付きの組を指すものとし，無限集合としての基底は正式な定義の対象としない．
\end{mycolumn}

基底のもっとも重要な性質は，「表し方がただ一通りに定まる」ことである．

\begin{theorem}{基底による表示の一意性}{基底による表示の一意性}
  $V$を体$K$上の線型空間とし，$\bm{v}_1, \bm{v}_2, \dots, \bm{v}_n \in V$とする．
  このとき，$(\bm{v}_1, \bm{v}_2, \dots, \bm{v}_n)$が$V$の基底であるための必要十分条件は，
  任意の$\bm{v} \in V$が
  \[
    \bm{v} = x_1 \bm{v}_1 + x_2 \bm{v}_2 + \cdots + x_n \bm{v}_n
    \quad (x_i \in K)
  \]
  の形にただ一通りに表されることである．
\end{theorem}

\begin{leftbar}
\begin{proof}
  まず，$(\bm{v}_1, \dots, \bm{v}_n)$が基底であるとする．
  生成性より，任意の$\bm{v} \in V$は上の形に表せる．
  いま，
  \[
    \bm{v} = \sum_{i=1}^{n} x_i \bm{v}_i = \sum_{i=1}^{n} y_i \bm{v}_i
  \]
  と二通りに表せたとすると，辺々引いて
  \[
    \sum_{i=1}^{n} (x_i - y_i) \bm{v}_i = \bm{0}_V
  \]
  を得る．線型独立性より$x_i - y_i = 0$，すなわち$x_i = y_i$である．
  よって表し方は一通りである．

  逆に，任意の$\bm{v} \in V$がただ一通りに表されるとする．
  表せること自体から$V = \Span(\bm{v}_1, \dots, \bm{v}_n)$であり，生成性が成り立つ．
  また，
  \[
    \sum_{i=1}^{n} \lambda_i \bm{v}_i = \bm{0}_V
  \]
  とすると，$\bm{0}_V$の表し方としては$\lambda_1 = \cdots = \lambda_n = 0$とするものがすでにある．
  一意性より$\lambda_i = 0$であり，線型独立性が成り立つ．
  以上の考察により，定理が示された．
\end{proof}
\end{leftbar}

\thref{基底による表示の一意性}により，基底を一つ固定すると，$V$のベクトルに$n$個の数の組が一意に対応する．

\begin{definition}{座標ベクトル}{座標ベクトル}
  $\mathcal{B} = (\bm{v}_1, \bm{v}_2, \dots, \bm{v}_n)$を線型空間$V$の基底とする．
  $\bm{v} \in V$に対して，\thref{基底による表示の一意性}により一意に定まる$x_1, x_2, \dots, x_n \in K$を用いて
  \[
    [\bm{v}]_{\mathcal{B}}
    = \begin{pmatrix} x_1 \\ x_2 \\ \vdots \\ x_n \end{pmatrix} \in K^n
  \]
  と定め，これを基底$\mathcal{B}$に関する$\bm{v}$の\textbf{座標ベクトル}\index{ざひょうべくとる@座標ベクトル}という．
\end{definition}

\begin{prop}{座標ベクトルの性質}{座標ベクトルの性質}
  $\mathcal{B}$を線型空間$V$の基底とする．
  このとき，写像
  \[
    V \to K^n, \quad \bm{v} \mapsto [\bm{v}]_{\mathcal{B}}
  \]
  は全単射であり，任意の$\bm{v}, \bm{w} \in V$および任意の$\alpha \in K$に対して
  \[
    [\bm{v} + \bm{w}]_{\mathcal{B}} = [\bm{v}]_{\mathcal{B}} + [\bm{w}]_{\mathcal{B}},
    \quad
    [\alpha \bm{v}]_{\mathcal{B}} = \alpha [\bm{v}]_{\mathcal{B}}
  \]
  が成り立つ．
\end{prop}

\begin{leftbar}
\begin{proof}
  $\mathcal{B} = (\bm{v}_1, \dots, \bm{v}_n)$とする．
  全射性は生成性から，単射性は表示の一意性から従う．

  $[\bm{v}]_{\mathcal{B}} = (x_i)$，$[\bm{w}]_{\mathcal{B}} = (y_i)$とすると，
  \[
    \bm{v} + \bm{w}
    = \sum_{i=1}^{n} x_i \bm{v}_i + \sum_{i=1}^{n} y_i \bm{v}_i
    = \sum_{i=1}^{n} (x_i + y_i) \bm{v}_i
  \]
  であり，表示の一意性より$[\bm{v} + \bm{w}]_{\mathcal{B}} = (x_i + y_i)$である．
  スカラー倍についても同様である．
  これが証明すべきことであった．
\end{proof}
\end{leftbar}

\begin{remark}[基底による$K^n$との同一視]
  \prref{座標ベクトルの性質}は，「基底を一つ選べば，抽象的な線型空間$V$は数ベクトル空間$K^n$と同じものとみなせる」ことを述べている．
  たとえば，次数$2$以下の実数係数多項式全体は，基底$(1, x, x^2)$を選ぶことで$\mathbb{R}^3$と同一視され，
  \[
    a + bx + cx^2 \longleftrightarrow \begin{pmatrix} a \\ b \\ c \end{pmatrix}
  \]
  という対応のもとで，多項式の計算が数ベクトルの計算に翻訳される．

  ただし，この同一視は基底の選び方に依存する．
  「どの基底を選んでも変わらない性質」と「基底に依存する表示」を区別することが，以下の議論の主題である．
\end{remark}

\phantomsection
\subsection{次元}

基底の取り方は一通りではないが，「基底に含まれるベクトルの個数」は基底の取り方によらない．
これを示すのが本項の目標である．

\partref{線型独立性と階数}では，数ベクトル空間$K^m$に限って「線型独立な組の拡大」と「取替補題」を示した（\lemref{線型独立な組の拡大}，\lemref{取替補題}）．
本項では，同じ議論が一般の線型空間においてもそのまま通用することを確認する．
これは単なる繰り返しではなく，\partref{線型独立性と階数}の具体的な議論を抽象的な枠組みへ再構成する作業である．

まず，準備として次の補題を用意する．

\begin{lemma}{線型独立な組の拡大}{一般の線型独立な組の拡大}
  $\bm{v}_1, \bm{v}_2, \dots, \bm{v}_k$を線型空間$V$の線型独立なベクトルとし，
  \[
    \bm{w} \notin \Span(\bm{v}_1, \bm{v}_2, \dots, \bm{v}_k)
  \]
  とする．このとき，$\bm{v}_1, \bm{v}_2, \dots, \bm{v}_k, \bm{w}$は線型独立である．
\end{lemma}

\begin{leftbar}
\begin{proof}
  \[
    \lambda_1 \bm{v}_1 + \cdots + \lambda_k \bm{v}_k + \mu \bm{w} = \bm{0}_V
  \]
  とする．もし$\mu \ne 0$ならば，
  \[
    \bm{w} = \sum_{i=1}^{k} \left(-\frac{\lambda_i}{\mu}\right) \bm{v}_i
    \in \Span(\bm{v}_1, \dots, \bm{v}_k)
  \]
  となり仮定に反する．よって$\mu = 0$である．
  このとき$\sum_{i=1}^{k} \lambda_i \bm{v}_i = \bm{0}_V$となるので，$\bm{v}_1, \dots, \bm{v}_k$の線型独立性より$\lambda_1 = \cdots = \lambda_k = 0$である．
  これが証明すべきことであった．
\end{proof}
\end{leftbar}

次が，次元の概念を支える中心的な補題である．

\begin{lemma}{取替補題}{一般の取替補題}
  $V$を線型空間とし，$\bm{v}_1, \bm{v}_2, \dots, \bm{v}_n \in V$とする．
  $\bm{w}_1, \bm{w}_2, \dots, \bm{w}_m \in \Span(\bm{v}_1, \bm{v}_2, \dots, \bm{v}_n)$が線型独立ならば，
  \[
    m \le n
  \]
  が成り立つ．
\end{lemma}

\begin{leftbar}
\begin{proof}
  $W = \Span(\bm{v}_1, \dots, \bm{v}_n)$とおく．
  $k = 0, 1, \dots, m$に対して，次の主張$(\ast_k)$を$k$についての帰納法で示す．

  \begin{quote}
    $(\ast_k)$：$k \le n$であり，かつ$\bm{v}_1, \dots, \bm{v}_n$の番号を適当に付け替えることで
    \[
      W = \Span(\bm{w}_1, \dots, \bm{w}_k, \bm{v}_{k+1}, \dots, \bm{v}_n)
    \]
    とできる．
  \end{quote}

  $k = 0$のときは$W$の定義そのものであり，成り立つ．

  $k < m$として$(\ast_k)$を仮定する．
  $\bm{w}_{k+1} \in W$であるから，帰納法の仮定より
  \begin{equation}
    \bm{w}_{k+1}
    = \sum_{i=1}^{k} a_i \bm{w}_i + \sum_{j=k+1}^{n} b_j \bm{v}_j
    \label{eq:一般の取替補題の展開}
  \end{equation}
  と表せる．

  ここで，$b_{k+1} = \cdots = b_n = 0$と仮定してみる（$k = n$のときは第2の和が空なので自動的にこの場合である）．
  このとき\eqref{eq:一般の取替補題の展開}は
  \[
    a_1 \bm{w}_1 + \cdots + a_k \bm{w}_k + (-1)\bm{w}_{k+1} = \bm{0}_V
  \]
  と書き直せる．$\bm{w}_{k+1}$の係数は$-1 \ne 0$であるから，これは$\bm{w}_1, \dots, \bm{w}_{k+1}$の非自明な線型関係であり，線型独立性に矛盾する．

  よって$k < n$，すなわち$k + 1 \le n$であり，さらに$b_j \ne 0$となる$j$が存在する．
  残っている$\bm{v}_{k+1}, \dots, \bm{v}_n$の番号を付け替えて$b_{k+1} \ne 0$としてよい．
  このとき\eqref{eq:一般の取替補題の展開}を$\bm{v}_{k+1}$について解くと，
  \[
    \bm{v}_{k+1}
    = \frac{1}{b_{k+1}}\left( \bm{w}_{k+1} - \sum_{i=1}^{k} a_i \bm{w}_i - \sum_{j=k+2}^{n} b_j \bm{v}_j \right)
  \]
  となるので，
  \[
    \bm{v}_{k+1} \in \Span(\bm{w}_1, \dots, \bm{w}_{k+1}, \bm{v}_{k+2}, \dots, \bm{v}_n)
  \]
  である．
  したがって，右辺の線型包は$\bm{w}_1, \dots, \bm{w}_k, \bm{v}_{k+1}, \dots, \bm{v}_n$をすべて含む線型部分空間であり，\prref{線型包の最小性}より
  \[
    W = \Span(\bm{w}_1, \dots, \bm{w}_k, \bm{v}_{k+1}, \dots, \bm{v}_n)
    \subset \Span(\bm{w}_1, \dots, \bm{w}_{k+1}, \bm{v}_{k+2}, \dots, \bm{v}_n)
  \]
  を得る．逆の包含は明らかであるから，$(\ast_{k+1})$が成り立つ．

  以上より$(\ast_m)$が成り立ち，とくに$m \le n$である．
  これが証明すべきことであった．
\end{proof}
\end{leftbar}

\begin{theorem}{基底の個数の一意性}{基底の個数の一意性}
  $V$を線型空間とし，$(\bm{v}_1, \dots, \bm{v}_n)$と$(\bm{w}_1, \dots, \bm{w}_m)$がともに$V$の基底であるとする．
  このとき，$m = n$である．
\end{theorem}

\begin{leftbar}
\begin{proof}
  $(\bm{v}_1, \dots, \bm{v}_n)$の生成性より$\bm{w}_1, \dots, \bm{w}_m \in \Span(\bm{v}_1, \dots, \bm{v}_n)$であり，$\bm{w}_1, \dots, \bm{w}_m$は線型独立であるから，\lemref{一般の取替補題}より$m \le n$である．

  $(\bm{v}_1, \dots, \bm{v}_n)$と$(\bm{w}_1, \dots, \bm{w}_m)$の役割を入れ替えて同じ議論を行えば$n \le m$を得る．
  よって$m = n$である．
  これが証明すべきことであった．
\end{proof}
\end{leftbar}

\thref{基底の個数の一意性}により，次の定義が意味をもつ．

\begin{definition}{次元}{次元}
  線型空間$V$が基底$(\bm{v}_1, \bm{v}_2, \dots, \bm{v}_n)$をもつとき，
  $n$を$V$の\textbf{次元}\index{じげん@次元}といい，
  \[
    \dim_K V = n
  \]
  と表す．体$K$が文脈から明らかなときは，単に$\dim V$とかく．

  とくに，$\dim \{\bm{0}_V\} = 0$である．
\end{definition}

\begin{remark}[「次元」が意味をもつために]
  「次元が$n$である」という主張は，
  \begin{itemize}
    \item 基底が存在すること（存在性）
    \item その個数が基底の取り方によらないこと（一意性）
  \end{itemize}
  の二つがそろって初めて意味をもつ．
  \thref{基底の個数の一意性}は後者を保証するものであり，前者は，有限生成な線型空間については次の\thref{基底の存在}が保証する（一般の線型空間の場合については，\deref{有限生成}のあとのコラムで触れた）．

  定義の順序としては，まず基底を定め，つぎにその個数が基底の取り方によらないことを示し，そのうえで次元を定義する，という流れになっている．次元を先に定義することはできない点に注意してほしい．
\end{remark}

次元に関連して，用語を一つ補っておく．

\begin{definition}{有限次元と無限次元}{有限次元と無限次元}
  有限個のベクトルからなる基底をもつ線型空間を\textbf{有限次元}\index{ゆうげんじげん@有限次元}線型空間という．
  有限次元でない線型空間を\textbf{無限次元}\index{むげんじげん@無限次元}線型空間という．
\end{definition}

のちの\thref{基底の存在}で示すように，有限生成な線型空間はつねに（有限個のベクトルからなる）基底をもつ．
逆に，有限個のベクトルからなる基底をもつ線型空間は，明らかに有限生成である．
したがって「有限生成」と「有限次元」は同じクラスの線型空間を指す言葉である．本書では，有限次元性（有限生成性）を仮定する命題では，その都度これを明記する（\deref{有限生成}のあとのコラムも参照）．

\begin{example}{次元の例}{次元の例}
  \exref{基底の例}より，
  \[
    \dim_K K^n = n,
    \quad
    \dim_K M_{m,n}(K) = mn,
    \quad
    \dim_{\mathbb{R}} \mathbb{R}[x]_{\le n} = n + 1
  \]
  である．ここで$M_{m,n}(K)$は$K$の元を成分とする$m \times n$行列全体の集合（\deref{行列全体の集合}），$\mathbb{R}[x]_{\le n}$は次数が$n$以下の実数係数多項式全体を表す．

  なお，$\mathbb{C}$は$\mathbb{C}$上の線型空間としては$\dim_{\mathbb{C}} \mathbb{C} = 1$であるが，$\mathbb{R}$上の線型空間とみなすと$(1, i)$が基底となり$\dim_{\mathbb{R}} \mathbb{C} = 2$である．
  次元は，どの体上の線型空間とみるかに依存する．
\end{example}

\phantomsection
\subsection{基底の存在と延長}

\begin{theorem}{基底の存在}{基底の存在}
  $V$を有限生成な線型空間とすると，$V$は基底をもつ．
  より詳しく，$V = \Span(\bm{v}_1, \dots, \bm{v}_n)$ならば，$\bm{v}_1, \dots, \bm{v}_n$のうちの適当な部分の組が$V$の基底となる．
\end{theorem}

\begin{leftbar}
\begin{proof}
  $\bm{v}_1, \dots, \bm{v}_n$が線型独立ならば，これがそのまま基底である．

  線型独立でないとすると，\prref{線型従属の言い換え}より，ある$i$について$\bm{v}_i$は残りのベクトルの線型結合として表せる．
  このとき，
  \[
    \Span(\bm{v}_1, \dots, \bm{v}_n)
    = \Span(\bm{v}_1, \dots, \widehat{\bm{v}_i}, \dots, \bm{v}_n)
  \]
  が成り立つ（$\widehat{\phantom{x}}$はその項を除くことを表す）．
  実際，右辺は左辺に含まれ，逆に$\bm{v}_i$が右辺に属するので，\prref{線型包の最小性}より左辺は右辺に含まれる．

  この操作でベクトルの個数は$1$減るので，有限回で線型独立な組に到達する．
  すべてのベクトルが除かれた場合は$V = \Span(\varnothing) = \{\bm{0}_V\}$であり，空の組が基底である．
  これが証明すべきことであった．
\end{proof}
\end{leftbar}

\begin{theorem}{基底の延長}{基底の延長}
  $V$を有限生成な線型空間とし，$\bm{w}_1, \bm{w}_2, \dots, \bm{w}_k \in V$を線型独立なベクトルとする．
  このとき，適当な$\bm{u}_1, \dots, \bm{u}_l \in V$を付け加えて
  \[
    (\bm{w}_1, \dots, \bm{w}_k, \bm{u}_1, \dots, \bm{u}_l)
  \]
  を$V$の基底とすることができる．
\end{theorem}

\begin{leftbar}
\begin{proof}
  $V = \Span(\bm{v}_1, \dots, \bm{v}_n)$とする．

  $\bm{v}_1, \bm{v}_2, \dots, \bm{v}_n$を順に見ていき，
  「その時点までにできている組の線型包に属さないもの」だけを組に付け加える，という操作を行う．
  \lemref{一般の線型独立な組の拡大}より，付け加えるたびに組は線型独立なままである．
  最終的に得られる組を
  \[
    (\bm{w}_1, \dots, \bm{w}_k, \bm{u}_1, \dots, \bm{u}_l)
  \]
  とおく．

  構成法より，各$\bm{v}_j$はこの組の線型包に属する（付け加えられた場合は自明であり，付け加えられなかった場合はその時点の線型包に属していたからである）．
  よって\prref{線型包の最小性}より
  \[
    V = \Span(\bm{v}_1, \dots, \bm{v}_n)
    \subset \Span(\bm{w}_1, \dots, \bm{w}_k, \bm{u}_1, \dots, \bm{u}_l)
    \subset V
  \]
  となり，生成性が成り立つ．線型独立性はすでに示したので，この組は$V$の基底である．
  これが証明すべきことであった．
\end{proof}
\end{leftbar}

\begin{cor}{次元による基底の判定}{次元による基底の判定}
  $\dim V = n$とし，$\bm{v}_1, \bm{v}_2, \dots, \bm{v}_n \in V$とする（個数が$n$であることに注意）．
  このとき，次の三つは同値である．
  \begin{enumerate}[label=(\roman*),leftmargin=3.2em]
    \item $\bm{v}_1, \dots, \bm{v}_n$は線型独立である．
    \item $\Span(\bm{v}_1, \dots, \bm{v}_n) = V$である．
    \item $(\bm{v}_1, \dots, \bm{v}_n)$は$V$の基底である．
  \end{enumerate}
\end{cor}

\begin{leftbar}
\begin{proof}
  (iii)から(i)，(ii)が従うのは基底の定義そのものである．

  (i)$\limp$(iii)：\thref{基底の延長}より，この組はベクトルを付け加えて基底にできる．
  しかし\thref{基底の個数の一意性}より基底の個数は$n$でなければならないから，付け加えるベクトルは$0$個である．

  (ii)$\limp$(iii)：\thref{基底の存在}より，この組から適当な部分を取り出して基底にできる．
  同様に，基底の個数は$n$であるから，取り除くベクトルは$0$個である．
  これが証明すべきことであった．
\end{proof}
\end{leftbar}

\begin{remark}[次元による基底の判定の注意]
  \coref{次元による基底の判定}は実用上たいへん便利である．
  たとえば$\mathbb{R}^3$のベクトルが$3$個与えられたとき，基底であることを確かめるには線型独立性だけを調べればよく，生成性を別に確認する必要はない．

  ただし，これは「個数が次元と一致している」場合に限られる．
  個数が$n$と異なる組については，線型独立性と生成性は独立な条件であり，一方から他方は従わない．
\end{remark}

\begin{prop}{部分空間の次元}{部分空間の次元}
  $V$を有限生成な線型空間とし，$W$を$V$の線型部分空間とする．
  このとき$W$も有限生成であり，
  \[
    \dim W \le \dim V
  \]
  が成り立つ．さらに，$\dim W = \dim V$であることと$W = V$であることは同値である．
\end{prop}

\begin{leftbar}
\begin{proof}
  $\dim V = n$とする．
  \lemref{一般の取替補題}より，$W$の線型独立な組に含まれるベクトルの個数は高々$n$である．
  空の組も線型独立とみなすと，$W$の線型独立な組の長さは$0, 1, \dots, n$のいずれかであるから，その最大値$r$が存在する．
  そこで，長さが最大である$W$の線型独立な組を$(\bm{w}_1, \dots, \bm{w}_r)$とする（$r \le n$．$W = \{\bm{0}_V\}$のときは$r = 0$であり，空の組をとることになる）．

  もし$\Span(\bm{w}_1, \dots, \bm{w}_r) \ne W$ならば，$\bm{w} \in W$であって$\bm{w} \notin \Span(\bm{w}_1, \dots, \bm{w}_r)$となるものがとれる．
  \lemref{一般の線型独立な組の拡大}より$\bm{w}_1, \dots, \bm{w}_r, \bm{w}$は線型独立となり，$r$の最大性に矛盾する．
  よって$(\bm{w}_1, \dots, \bm{w}_r)$は$W$の基底であり，$\dim W = r \le n$である．

  最後の主張について，$W = V$ならば$\dim W = \dim V$は明らかである．
  逆に$\dim W = n$とすると，$W$の基底は$V$の$n$個の線型独立なベクトルであるから，\coref{次元による基底の判定}より$V$の基底でもある．
  よって
  \[
    V = \Span(\bm{w}_1, \dots, \bm{w}_n) \subset W \subset V
  \]
  となり，$W = V$である．
  これが証明すべきことであった．
\end{proof}
\end{leftbar}

\begin{theorem}{次元公式}{次元公式}
  $V$を有限生成な線型空間とし，$W_1, W_2$を$V$の線型部分空間とする．
  このとき，
  \[
    \dim(W_1 + W_2) + \dim(W_1 \cap W_2) = \dim W_1 + \dim W_2
  \]
  が成り立つ．
\end{theorem}

\begin{leftbar}
\begin{proof}
  $\dim(W_1 \cap W_2) = r$とし，$W_1 \cap W_2$の基底を$(\bm{u}_1, \dots, \bm{u}_r)$とする．
  \thref{基底の延長}より，これを延長して
  \[
    W_1 \text{の基底} \quad (\bm{u}_1, \dots, \bm{u}_r, \bm{x}_1, \dots, \bm{x}_s),
  \]
  \[
    W_2 \text{の基底} \quad (\bm{u}_1, \dots, \bm{u}_r, \bm{y}_1, \dots, \bm{y}_t)
  \]
  がとれる．すなわち$\dim W_1 = r + s$，$\dim W_2 = r + t$である．

  ここで，
  \[
    (\bm{u}_1, \dots, \bm{u}_r, \bm{x}_1, \dots, \bm{x}_s, \bm{y}_1, \dots, \bm{y}_t)
  \]
  が$W_1 + W_2$の基底であることを示す．

  生成性について，まず，上の組の各ベクトルは$W_1$または$W_2$に属するから$W_1 + W_2$に属し，$W_1 + W_2$は線型部分空間であるから，\prref{線型包の最小性}より上の組の線型包は$W_1 + W_2$に含まれる．
  逆に，$W_1$も$W_2$も上の組の線型包に含まれるから，\prref{和と線型包}の$W_1 + W_2 = \Span(W_1 \cup W_2)$と\prref{線型包の最小性}より，$W_1 + W_2$は上の組の線型包に含まれる．
  したがって，上の組の線型包は$W_1 + W_2$に等しい．

  線型独立性について，
  \begin{equation}
    \sum_{i=1}^{r} a_i \bm{u}_i + \sum_{j=1}^{s} b_j \bm{x}_j + \sum_{k=1}^{t} c_k \bm{y}_k = \bm{0}_V
    \label{eq:次元公式の線型関係}
  \end{equation}
  とする．このとき
  \[
    \bm{z} := \sum_{k=1}^{t} c_k \bm{y}_k
    = -\left( \sum_{i=1}^{r} a_i \bm{u}_i + \sum_{j=1}^{s} b_j \bm{x}_j \right)
  \]
  とおくと，左辺の表示より$\bm{z} \in W_2$，右辺の表示より$\bm{z} \in W_1$であるから，$\bm{z} \in W_1 \cap W_2$である．
  よって
  \[
    \bm{z} = \sum_{i=1}^{r} d_i \bm{u}_i
  \]
  と表せるので，
  \[
    \sum_{i=1}^{r} d_i \bm{u}_i - \sum_{k=1}^{t} c_k \bm{y}_k = \bm{0}_V
  \]
  となる．$(\bm{u}_1, \dots, \bm{u}_r, \bm{y}_1, \dots, \bm{y}_t)$は$W_2$の基底であるから線型独立であり，とくに$c_1 = \cdots = c_t = 0$である．

  これを\eqref{eq:次元公式の線型関係}に代入すると
  \[
    \sum_{i=1}^{r} a_i \bm{u}_i + \sum_{j=1}^{s} b_j \bm{x}_j = \bm{0}_V
  \]
  となり，$W_1$の基底の線型独立性より$a_i = 0$，$b_j = 0$である．

  以上より$\dim(W_1 + W_2) = r + s + t$であり，
  \[
    (r + s + t) + r = (r + s) + (r + t)
  \]
  から主張が従う．
  これが証明すべきことであった．
\end{proof}
\end{leftbar}

\begin{remark}[次元公式の読み方]
  次元公式は，有限集合の元の個数についての公式
  \[
    |A \cup B| + |A \cap B| = |A| + |B|
  \]
  と同じ形をしている．
  ただし，線型空間では合併$W_1 \cup W_2$ではなく和$W_1 + W_2$を考える点が異なる（\prref{和と線型包}のあとのRemarkを参照）．

  とくに$W_1 \cap W_2 = \{\bm{0}_V\}$のときは
  \[
    \dim(W_1 + W_2) = \dim W_1 + \dim W_2
  \]
  となる．この場合，和$W_1 + W_2$は（それ自身を一つの線型空間とみて）$W_1$と$W_2$の内部直和$W_1 \oplus W_2$である（\deref{内部直和}）．
  このとき各元が$\bm{w}_1 + \bm{w}_2$（$\bm{w}_i \in W_i$）の形にただ一通りに表されることは，\prref{内部直和と分解の一意性}で見た通りである．
  すなわち上の等式は，「内部直和の次元は各成分の次元の和である」と読むことができる．
\end{remark}

本項の締めくくりとして，もう一つの次元公式を証明しよう．
\partref{線型写像}で，線型写像$f \colon K^n \to K^m$に対して像$\Im f$と核$\Ker f$を定義した（\deref{像と核}）．
本章では，これらがそれぞれ$K^m$，$K^n$の線型部分空間であることを確かめた（\prref{像と核の線型部分空間性}）ので，その次元を考えることができる．

\begin{theorem}{線型写像の次元公式}{線型写像の次元公式}
  $f \colon K^n \to K^m$を線型写像とする．このとき，
  \[
    \dim \Ker f + \dim \Im f = n
  \]
  が成り立つ．
\end{theorem}

\begin{leftbar}
\begin{proof}
  $K^n$は有限生成であり（標準基底，\exref{基底の例}），$\Ker f$はその線型部分空間であるから，\prref{部分空間の次元}より$\Ker f$も有限生成である．
  $\dim \Ker f = k$とおき，$\Ker f$の基底を$(\bm{u}_1, \dots, \bm{u}_k)$とする（$k = 0$のときは空の組をとる）．
  \thref{基底の延長}より，これを延長して$K^n$の基底
  \[
    (\bm{u}_1, \dots, \bm{u}_k, \bm{w}_1, \dots, \bm{w}_l)
  \]
  がとれる．$\dim K^n = n$（\exref{次元の例}）であるから，\thref{基底の個数の一意性}より$k + l = n$である．

  ここで，$(f(\bm{w}_1), \dots, f(\bm{w}_l))$が$\Im f$の基底であることを示す．

  生成性について，任意の$\bm{v} \in \Im f$は，ある$\bm{x} \in K^n$を用いて$\bm{v} = f(\bm{x})$と表せる．
  $\bm{x}$を上の$K^n$の基底で
  \[
    \bm{x} = \sum_{i=1}^{k} a_i \bm{u}_i + \sum_{j=1}^{l} b_j \bm{w}_j
  \]
  と表すと，\prref{線型写像と線型結合}と$f(\bm{u}_i) = \bm{0}$（$\bm{u}_i \in \Ker f$）より，
  \[
    \bm{v} = f(\bm{x})
    = \sum_{i=1}^{k} a_i f(\bm{u}_i) + \sum_{j=1}^{l} b_j f(\bm{w}_j)
    = \sum_{j=1}^{l} b_j f(\bm{w}_j)
  \]
  となる．よって$\Im f = \Span\bigl(f(\bm{w}_1), \dots, f(\bm{w}_l)\bigr)$である．

  線型独立性について，
  \[
    \sum_{j=1}^{l} c_j f(\bm{w}_j) = \bm{0}
  \]
  とする．\prref{線型写像と線型結合}より$f\left(\sum_{j=1}^{l} c_j \bm{w}_j\right) = \bm{0}$，すなわち
  \[
    \sum_{j=1}^{l} c_j \bm{w}_j \in \Ker f
  \]
  である．$(\bm{u}_1, \dots, \bm{u}_k)$は$\Ker f$の基底であるから，
  \[
    \sum_{j=1}^{l} c_j \bm{w}_j = \sum_{i=1}^{k} d_i \bm{u}_i
  \]
  と表せる．移項すると
  \[
    \sum_{i=1}^{k} d_i \bm{u}_i - \sum_{j=1}^{l} c_j \bm{w}_j = \bm{0}
  \]
  となり，$(\bm{u}_1, \dots, \bm{u}_k, \bm{w}_1, \dots, \bm{w}_l)$の線型独立性より，とくに$c_1 = \cdots = c_l = 0$である．

  以上より$(f(\bm{w}_1), \dots, f(\bm{w}_l))$は$\Im f$の基底であり，$\dim \Im f = l$である．
  したがって，
  \[
    \dim \Ker f + \dim \Im f = k + l = n
  \]
  が成り立つ．これが証明すべきことであった．
\end{proof}
\end{leftbar}

\partref{線型独立性と階数}で「線型独立な列の最大個数」として導入した階数（\deref{階数}）は，この公式を通じて「像の次元」として捉え直される．

\begin{cor}{階数と像・核の次元}{階数と像の次元}
  $A$を$m \times n$行列とし，$T_A \colon K^n \to K^m$を$A$の定める線型写像とする．このとき，
  \[
    \dim \Im T_A = \rank A,
    \qquad
    \dim \Ker T_A = n - \rank A
  \]
  が成り立つ．
\end{cor}

\begin{leftbar}
\begin{proof}
  $A$の列ベクトルを$\bm{a}_1, \dots, \bm{a}_n \in K^m$とする．
  任意の$\bm{x} \in K^n$に対して，$\bm{x} = \sum_{j=1}^{n} x_j \bm{e}_j$，$A\bm{e}_j = \bm{a}_j$と\prref{線型写像と線型結合}より
  \[
    A\bm{x} = \sum_{j=1}^{n} x_j \bm{a}_j
  \]
  であるから，$\Im T_A = \Span(\bm{a}_1, \dots, \bm{a}_n)$である．

  $r = \rank A$とおく．
  階数の定義（\deref{階数}）より，$\bm{a}_1, \dots, \bm{a}_n$のうちから線型独立な$r$本$\bm{a}_{j_1}, \dots, \bm{a}_{j_r}$が選べる．
  任意の列$\bm{a}_j$に対して，$r + 1$本の組$\bm{a}_{j_1}, \dots, \bm{a}_{j_r}, \bm{a}_j$は$r$の最大性より線型従属であるから，\lemref{一般の線型独立な組の拡大}の対偶より
  \[
    \bm{a}_j \in \Span(\bm{a}_{j_1}, \dots, \bm{a}_{j_r})
  \]
  である．よって\prref{線型包の最小性}より
  \[
    \Im T_A = \Span(\bm{a}_1, \dots, \bm{a}_n)
    \subset \Span(\bm{a}_{j_1}, \dots, \bm{a}_{j_r})
    \subset \Im T_A
  \]
  となり，$\Im T_A = \Span(\bm{a}_{j_1}, \dots, \bm{a}_{j_r})$である．
  $(\bm{a}_{j_1}, \dots, \bm{a}_{j_r})$は線型独立かつ$\Im T_A$を生成するから，$\Im T_A$の基底であり，$\dim \Im T_A = r = \rank A$である．

  後半の等式は，\thref{線型写像の次元公式}より
  \[
    \dim \Ker T_A = n - \dim \Im T_A = n - \rank A
  \]
  として従う．これが証明すべきことであった．
\end{proof}
\end{leftbar}

\begin{remark}[次元公式の解釈と解の自由度]
  \thref{線型写像の次元公式}は，定義域のもつ$n$次元分の情報が，「核としてつぶれる分」と「像として生き残る分」とに過不足なく分かれることを述べている．
  \exref{像と核の例}の正射影$f_3$でいえば，$\dim \Ker f_3 = 1$（$y$軸方向がつぶれる）と$\dim \Im f_3 = 1$（$x$軸が残る）の和が，定義域$\mathbb{R}^2$の次元$2$に一致している．
  文献によっては，\thref{次元公式}ではなくこちらを指して「次元公式」あるいは「次元定理」とよぶことも多い．

  また，\coref{階数と像の次元}の$\dim \Ker T_A = n - \rank A$は，連立一次方程式の解の様子を定量的に述べている．
  斉次連立一次方程式$A\bm{x} = \bm{0}$の解空間は$\Ker T_A$であり（\prref{連立一次方程式と像と核}），その次元$n - \rank A$は，解を表すのに必要なパラメータの個数（解の自由度）にあたる．
  とくに$\rank A = n$のときは$\Ker T_A = \{\bm{0}\}$となり，\prref{単射性と核}より$T_A$は単射，すなわち$A\bm{x} = \bm{0}$は自明な解$\bm{x} = \bm{0}$しかもたない．

  なお，のちに一般の線型空間のあいだの線型写像を導入すれば，有限次元の線型空間$V$から線型空間$W$への線型写像$f$に対しても，まったく同じ証明によって
  \[
    \dim \Ker f + \dim \Im f = \dim V
  \]
  が成り立つ．
\end{remark}

\phantomsection
\section{基底の取替え}

ベクトルそのものは，どの基底を選ぶかとは無関係に定まっている．
一方，その座標ベクトルは，基底の選び方に応じて変化する「表示」である（\prref{座標ベクトルの性質}のあとのRemarkを参照）．
本節では，この「対象」と「表示」の区別を踏まえ，基底を取り替えたときに座標ベクトルがどのような規則で変換されるかを調べる．

\phantomsection
\subsection{基底の行列表記}

基底を取り替える操作を見通しよく扱うために，記法を一つ用意する．

\begin{definition}{ベクトルの組と行列の積}{ベクトルの組と行列の積}
  $V$を体$K$上の線型空間とし，$\bm{v}_1, \bm{v}_2, \dots, \bm{v}_n \in V$とする．
  $n \times m$行列$P = (p_{ij})$に対して，
  \[
    (\bm{v}_1, \bm{v}_2, \dots, \bm{v}_n) P
    :=
    \left( \sum_{i=1}^{n} p_{i1} \bm{v}_i,~ \sum_{i=1}^{n} p_{i2} \bm{v}_i,~ \dots,~ \sum_{i=1}^{n} p_{im} \bm{v}_i \right)
  \]
  と定める．すなわち，右辺の第$j$成分は$P$の第$j$列を係数とする$\bm{v}_1, \dots, \bm{v}_n$の線型結合である．

  同様に，$\bm{x} = (x_i) \in K^n$に対して
  \[
    (\bm{v}_1, \bm{v}_2, \dots, \bm{v}_n) \bm{x}
    := \sum_{i=1}^{n} x_i \bm{v}_i \in V
  \]
  と定める．
\end{definition}

\begin{remark}[形式的な行列記法について]
  これは，ベクトルを成分とする「$1 \times n$行列」と通常の行列との積を，形式的に行列の積の定義（\deref{行列の積}）と同じ規則で定めたものである．
  $V$が一般の線型空間の場合には成分どうしの積は定義されないので，あくまで記法の約束として理解してほしい．

  この記法のもとで，基底$\mathcal{B} = (\bm{v}_1, \dots, \bm{v}_n)$と座標ベクトルの関係は
  \[
    \bm{v} = \mathcal{B} \, [\bm{v}]_{\mathcal{B}}
  \]
  と簡潔にかける．
\end{remark}

\begin{prop}{ベクトルの組と行列の積の結合律}{ベクトルの組と行列の積の結合律}
  $\bm{v}_1, \dots, \bm{v}_n \in V$とし，$P$を$n \times m$行列，$Q$を$m \times l$行列とする．
  $\mathcal{V} = (\bm{v}_1, \dots, \bm{v}_n)$とかくと，
  \[
    (\mathcal{V} P) Q = \mathcal{V} (PQ)
  \]
  が成り立つ．
\end{prop}

\begin{leftbar}
\begin{proof}
  $P = (p_{ij})$，$Q = (q_{jk})$とし，$\mathcal{V}P$の第$j$成分を$\bm{w}_j = \sum_{i=1}^{n} p_{ij} \bm{v}_i$とおく．
  このとき，$(\mathcal{V}P)Q$の第$k$成分は
  \begin{align*}
    \sum_{j=1}^{m} q_{jk} \bm{w}_j
    &= \sum_{j=1}^{m} q_{jk} \left( \sum_{i=1}^{n} p_{ij} \bm{v}_i \right) \\
    &= \sum_{i=1}^{n} \left( \sum_{j=1}^{m} p_{ij} q_{jk} \right) \bm{v}_i
  \end{align*}
  である．最後の括弧の中は$PQ$の$(i,k)$成分にほかならないから，これは$\mathcal{V}(PQ)$の第$k$成分に等しい．
  これが証明すべきことであった．
\end{proof}
\end{leftbar}

\phantomsection
\subsection{基底の取替行列}

\begin{definition}{基底の取替行列}{取替行列}
  $\mathcal{B} = (\bm{v}_1, \dots, \bm{v}_n)$，$\mathcal{B}' = (\bm{v}'_1, \dots, \bm{v}'_n)$をともに線型空間$V$の基底とする．
  各$\bm{v}'_j$を$\mathcal{B}$で表したときの係数を$p_{ij}$とする．すなわち，
  \[
    \bm{v}'_j = \sum_{i=1}^{n} p_{ij} \bm{v}_i \quad (1 \le j \le n)
  \]
  とし，$P = (p_{ij})$とおく．
  このとき，
  \[
    \mathcal{B}' = \mathcal{B} P
  \]
  が成り立つ．この$P$を，$\mathcal{B}$から$\mathcal{B}'$への\textbf{基底の取替行列}\index{きていのとりかええぎょうれつ@基底の取替行列}という．
\end{definition}

\begin{remark}[取替行列の列の意味]
  取替行列$P$の第$j$列は，新しい基底ベクトル$\bm{v}'_j$の，古い基底$\mathcal{B}$に関する座標ベクトルである：
  \[
    P = \bigl( [\bm{v}'_1]_{\mathcal{B}},~ [\bm{v}'_2]_{\mathcal{B}},~ \dots,~ [\bm{v}'_n]_{\mathcal{B}} \bigr)
  \]
  「新しい基底を，古い基底で表したものを列に並べる」と覚えるとよい．
  添字の付き方が$\bm{v}'_j = \sum_i p_{ij} \bm{v}_i$と，和をとるのが第$1$添字である点に注意してほしい．
  これは$\mathcal{B}' = \mathcal{B}P$という表記が成り立つように定めた結果である．
\end{remark}

\begin{prop}{取替行列の正則性}{取替行列の正則性}
  $\mathcal{B}$から$\mathcal{B}'$への取替行列を$P$とする．
  このとき$P$は正則であり，$\mathcal{B}'$から$\mathcal{B}$への取替行列は$P^{-1}$である．
\end{prop}

\begin{leftbar}
\begin{proof}
  $\mathcal{B}'$から$\mathcal{B}$への取替行列を$Q$とすると，
  \[
    \mathcal{B} = \mathcal{B}' Q = (\mathcal{B}P)Q = \mathcal{B}(PQ)
  \]
  である（最後の等号は\prref{ベクトルの組と行列の積の結合律}による）．
  一方，明らかに$\mathcal{B} = \mathcal{B}E_n$である．
  $\mathcal{B}$は基底であるから表示は一意であり（\thref{基底による表示の一意性}），両辺の各成分を比較して
  \[
    PQ = E_n
  \]
  を得る．$\mathcal{B}$と$\mathcal{B}'$の役割を入れ替えれば$QP = E_n$も同様に従う．
  よって$P$は正則で$Q = P^{-1}$である．
  これが証明すべきことであった．
\end{proof}
\end{leftbar}

\phantomsection
\subsection{座標の変換}

基底を取り替えると，同じベクトルの座標ベクトルはどう変わるだろうか．

\begin{theorem}{座標の変換公式}{座標の変換公式}
  $\mathcal{B}$，$\mathcal{B}'$を線型空間$V$の基底とし，$P$を$\mathcal{B}$から$\mathcal{B}'$への取替行列とする．
  このとき，任意の$\bm{v} \in V$に対して
  \[
    [\bm{v}]_{\mathcal{B}} = P \, [\bm{v}]_{\mathcal{B}'},
    \quad \text{すなわち} \quad
    [\bm{v}]_{\mathcal{B}'} = P^{-1} [\bm{v}]_{\mathcal{B}}
  \]
  が成り立つ．
\end{theorem}

\begin{leftbar}
\begin{proof}
  $\bm{x} = [\bm{v}]_{\mathcal{B}}$，$\bm{x}' = [\bm{v}]_{\mathcal{B}'}$とおく．
  座標ベクトルの定義より
  \[
    \bm{v} = \mathcal{B}\bm{x} = \mathcal{B}'\bm{x}'
  \]
  である．ここで$\mathcal{B}' = \mathcal{B}P$であるから，\prref{ベクトルの組と行列の積の結合律}より
  \[
    \mathcal{B}'\bm{x}' = (\mathcal{B}P)\bm{x}' = \mathcal{B}(P\bm{x}')
  \]
  となる．よって
  \[
    \mathcal{B}\bm{x} = \mathcal{B}(P\bm{x}')
  \]
  であり，$\mathcal{B}$による表示の一意性（\thref{基底による表示の一意性}）から$\bm{x} = P\bm{x}'$を得る．
  \prref{取替行列の正則性}より$P$は正則であるから，$\bm{x}' = P^{-1}\bm{x}$である．
  これが証明すべきことであった．
\end{proof}
\end{leftbar}

\begin{remark}[基底と座標の変換の向き]
  ここで，向きが「逆になる」ことに注意してほしい．
  \[
    \text{基底：} \quad \mathcal{B}' = \mathcal{B}P,
    \qquad
    \text{座標：} \quad [\bm{v}]_{\mathcal{B}'} = P^{-1} [\bm{v}]_{\mathcal{B}}
  \]
  基底は$P$によって新しいものに移るのに対し，座標は$P^{-1}$によって移るのである．
  この点は混乱しやすいので，具体例で確かめておこう．

  $\mathbb{R}^1$において$\bm{v}_1 = 1$を基底とし，$\bm{v}'_1 = 2$を新しい基底とする．
  このとき$\bm{v}'_1 = 2\bm{v}_1$であるから$P = (2)$である．
  ベクトル$\bm{v} = 6$の座標は，$\mathcal{B}$では$6$，$\mathcal{B}'$では$3$である．
  基底ベクトルを$2$倍に「伸ばす」と，同じベクトルを測る目盛りの数は$1/2$倍に「縮む」わけである．

  ものさしの目盛りを細かくすれば読み取る数値は大きくなり，粗くすれば小さくなる，という日常の感覚と同じである．
  このように，基底の変換と座標の変換とに互いに逆の行列（$P$と$P^{-1}$）が現れることを指して，座標は基底の変換に対して\textbf{反変的}\index{はんぺんてき@反変的}に変換する，と表現することがある．
\end{remark}

\begin{example}{座標の変換}{座標の変換}
  $V = \mathbb{R}^2$とし，$\mathcal{B} = (\bm{e}_1, \bm{e}_2)$を標準基底，
  \[
    \mathcal{B}' = (\bm{v}'_1, \bm{v}'_2), \quad
    \bm{v}'_1 = \begin{pmatrix} 1 \\ 1 \end{pmatrix}, \quad
    \bm{v}'_2 = \begin{pmatrix} 1 \\ -1 \end{pmatrix}
  \]
  とする．
  $\bm{v}'_1 = \bm{e}_1 + \bm{e}_2$，$\bm{v}'_2 = \bm{e}_1 - \bm{e}_2$であるから，$\mathcal{B}$から$\mathcal{B}'$への取替行列は
  \[
    P = \begin{pmatrix} 1 & 1 \\ 1 & -1 \end{pmatrix}
  \]
  である．$\det P = -2 \ne 0$より$P$は正則で，
  \[
    P^{-1} = \frac{1}{2}\begin{pmatrix} 1 & 1 \\ 1 & -1 \end{pmatrix}
  \]
  である．

  いま，$\bm{v} = \begin{pmatrix} 5 \\ 1 \end{pmatrix}$を考える．
  $\mathcal{B}$は標準基底であるから$[\bm{v}]_{\mathcal{B}} = \begin{pmatrix} 5 \\ 1 \end{pmatrix}$であり，
  \[
    [\bm{v}]_{\mathcal{B}'}
    = P^{-1}\begin{pmatrix} 5 \\ 1 \end{pmatrix}
    = \frac{1}{2}\begin{pmatrix} 6 \\ 4 \end{pmatrix}
    = \begin{pmatrix} 3 \\ 2 \end{pmatrix}
  \]
  となる．実際，
  \[
    3\begin{pmatrix} 1 \\ 1 \end{pmatrix} + 2\begin{pmatrix} 1 \\ -1 \end{pmatrix}
    = \begin{pmatrix} 5 \\ 1 \end{pmatrix}
  \]
  であり，確かに成り立っている．
\end{example}

\begin{remark}[標準基底の場合と行列式による判定]
  $V = K^n$で$\mathcal{B}$が標準基底の場合，取替行列$P$は「新しい基底ベクトルを列として並べた行列」そのものである．
  この場合，$\bm{v} = P[\bm{v}]_{\mathcal{B}'}$という関係は，\partref{線型独立性と階数}のRemarkで見た「$A\bm{x}$は$A$の列ベクトルの線型結合であり，その係数は$\bm{x}$の成分である」という事実にほかならない．

  また，$K = \mathbb{R}$または$\mathbb{C}$とする．\partref{行列式}では行列式に関する諸性質を複素数行列について証明したが，実数行列についても同じ定義式と証明がそのまま通用する（\partref{行列式}冒頭の注意を参照）．よって，\coref{次元による基底の判定}と\prref{行列式と階数}を合わせると，
  $K^n$の$n$個のベクトル$\bm{v}'_1, \dots, \bm{v}'_n$が基底をなすことは，それらを列に並べた行列$P$について$\det P \ne 0$が成り立つことと同値である．
  基底かどうかの判定が，行列式の計算に帰着されるのである．
\end{remark}
